\documentclass[11pt,a4paper]{article}

\usepackage[polish]{babel}
\usepackage[utf8]{inputenc}
\usepackage{polski}
\usepackage[T1]{fontenc}
\usepackage{indentfirst}
\usepackage{wrapfig}    % for wrapping figures, tables
\usepackage{isotope}
\usepackage{hyperref}
\usepackage{amsmath}
\usepackage{amssymb}
\usepackage{bm}
\usepackage{gensymb}
\usepackage{epsfig}
\usepackage{graphics}
\usepackage[shortlabels]{enumitem}
\usepackage{xspace}
\xspaceaddexceptions{[]\{\}}

%
%
%fixpagesize
\pagestyle{empty}
\addtolength{\textwidth}{4cm}
\addtolength{\textheight}{6cm}
\addtolength{\evensidemargin}{-3cm}
\addtolength{\oddsidemargin}{-2cm}
\addtolength{\topmargin}{-3cm}
\parindent=0cm

%
%
% Changes figure placing algorithm
\renewcommand{\topfraction}{1}       % maximal fraction of a page allowed for figures
\renewcommand{\textfraction}{0.15}   % minimal number of text for figure-text shared pages
\renewcommand{\floatpagefraction}{0.95} % if two above does not help, this could do the job 
                                        % must be: floatpagefraction < topfraction !!!!
\renewcommand{\textfraction}{0} % minimum fraction of page, which must be  devoted to text
\renewcommand{\topfraction}{1}  % maximum fraction at top, which can be occupied whit floats
\setcounter{totalnumber}{400}   % increase the number of floats for one page
\setcounter{topnumber}{200}     % at all/top/bottom.
\setcounter{bottomnumber}{200}  %

%
%
%small distance in list/item/enum for enumitem package
\setlist[itemize,enumerate]{topsep=0em}
\setlist{noitemsep}

%Nuclear notations: usage \nucl{235}{92}{U}. Math mode optional
\newcommand{\nucl}[3]{\ensuremath{
  \phantom{\ensuremath{^{#1}_{#2}}}
  \llap{\ensuremath{^{#1}}}
  \llap{\ensuremath{_{\rule{0pt}{.75em}#2}}}
  \mbox{#3} } 
} 

%print zadanie #
\newcounter{zadanie}\newcommand{\zadanie}[1][]{\addtocounter{zadanie}{1} ~\\  {\bf \emph{Zadanie \arabic{zadanie} #1 }} \\}
\newcounter{zaddom}\newcommand{\zaddom}[1][]{\addtocounter{zaddom}{1} ~\\  {\bf \emph{Zadanie domowe \arabic{zaddom} #1 }} \\}

%dynamic solutions hiding
\usepackage{comment}

\newif\ifshowsolutions

\def\showsolutions{} %comment this line to hide solutions

\newif\ifshowsolutions
\ifdefined\showsolutions
  \showsolutionstrue
\else
  \showsolutionsfalse
\fi

\ifshowsolutions
  \newenvironment{solution}{\vspace{1cm} \par\textbf{Rozwiązanie.} \vspace{1cm}}{}
\else
  \excludecomment{solution}
\fi
%%%%%%%%%%%%%%%%%%%%%%%%%%%%%%%%%%%%%%%%%%
\begin{document}   

%%%%%%%%%%%%%%%%%%%%%%%%%%%%%%%%%%%%%%%%%%%%%%%%%%%%%
\begin{centering}
  {\bf {\Large
  Termodynamika i Fizyka Statystyczna R \\
  Zestaw VI: Przekształcenie piekarskie.
  }} \\  
\end{centering}

\vspace{1cm}

%%%%%%%%%%%%%%%%%%%%%%%%%
\zadanie[]

Rozpatrzmy przekształcenie jednostkowego kwadratu: 
\begin{enumerate}
\item skalowanie o czynniki $2$ i $1/2$ w kierunkach $x$ i $y$, 
\item rozcięcie na pół w kierunku $x$ 
\item nałożenie jedną część nad drugą w kierunku $y$.
\end{enumerate}, czyli:
\begin{align*}
\vec{r}' = \hat{B}(\vec{r}) = \begin{cases} (2x, y/2), & x < 1/2 \\ (2x - 1, (y+1)/2), & x \geq 1/2 \end{cases}
\end{align*}

Proszę:
\begin{itemize}
\item znaleźć przekształcenie odwrotne: $\hat{B}^{-1}$, o ile takowe istnieje.
\item sprawdzić, czy przekształcenie piekarskie zachowuje pole w przestrzeni fazowej
\end{itemize}

\begin{solution}

Jesteśmy w stanie jednoznacznie odtworzyć punkt $(x, y)$ z jego obrazu $(x', y')$, 
więc przekształcenie jest odwracalne. Przekształcenie odwrotne ma postać:
\begin{align*}
(x, y) = \hat{B}^{-1}(x', y') = \begin{cases} (x'/2, 2y'), & y' < 1/2 \\ ((x'+1)/2, 2y' - 1), & y' \geq 1/2 \end{cases}
\end{align*}

Czy przekształcenie zachowuje pole w przestrzeni fazowej? 
Tak, wystarczy rozpatrzyć jakobian i pokazać, że jest równy $1$.
Zapiszmy przekształcenie odwrotne w postaci macierzowej:
\begin{align*}
\left( \begin{array}{c} x' \\ y' \end{array} \right) = \left( \begin{array}{cc} 2 & 0 \\
0 & 1/2 \end{array} \right) \left( \begin{array}{c} x \\ y \end{array} \right) + 
\left( \begin{array}{c} 0 \\ 0 \end{array} \right) \text{ lub } \left( \begin{array}{cc} 2 & 0 \\ 0 & 1/2 \end{array} \right) \left( \begin{array}{c} x \\ y \end{array} \right) + \left( \begin{array}{c} -1 \\ 1/2 \end{array} \right)
\end{align*}

Jakobian jest dany przez wyznacznik macierzy przekształcenia, 
która jest pochodną przekształcenia odwrotnego względem $x$ i $y$:
\begin{align*}
J = \det \left( \begin{array}{cc} 2 & 0 \\
0 & 1/2 \end{array} \right) = 2 \cdot \frac{1}{2} = 1
\end{align*}


\newpage
\end{solution}
%%%%%%%%%%%%%%%%%%%%%%%%%%%%%%%%%%%%%%%%%%%%%%%%%%%%%%%%%%%%%%%%%%%%%%%
%%%%%%%%%%%%%%%%%%%%%%%%%%%%%%%%%%%%%%%%%%%%%%%%%%%%%%%%%%%%%%%%%%%%%%%
\zadanie[]

Niech $\rho(x, y)$ opisuje gęstość prawdopodobieństwa znalezienia układu 
w danym punkcie $(x, y)$ przestrzeni fazowej. Proszę:
\begin{itemize}
    \item znaleźć postać gęstości prawdopodobieństwa po $n$ iteracjach 
    przekształcenia: $\rho_n(x, y)$
    \item dyskretny odpowiednik równania Liouville'a rządzącego ewolucją $\rho$ 
    podczas kolejnych iteracji przekształcenia: $\rho_{n+1} = f(\rho_n)$
\end{itemize}


\begin{solution}

Gęstość prawdopodobieństwa w kolejnej iteracji znajdujemy 
sumując gęstość prawdopodobieństwa w poprzedniej iteracji po wszystkich 
punktach,które są przekształcane do punktu $(x, y)$ w przestrzeni fazowej:
\begin{align*}
    \rho_{n+1}(x', y') = \int dx \int dy \rho_{n}(x, y) \delta((x', y') - \hat{B}(x, y))
\end{align*}

Odwzorowanie jest bijekcją, więc dla każdego punktu $(x', y')$ istnieje 
dokładnie jeden punkt $(x, y)$ taki że $\hat{B}(x, y) = (x', y')$. Oznacza to, że powyższy całkowy wzór można uprościć do:
\begin{align*}
    \rho_{n+1}(x', y') = \rho_n(\hat{B}^{-1}(x', y'))
\end{align*}

\newpage
\end{solution}
%%%%%%%%%%%%%%%%%%%%%%%%%%%%%%%%%%%%%%%%%%%%%%%%%%%%%%%%%%%%%%%%%%%%%%%
%%%%%%%%%%%%%%%%%%%%%%%%%%%%%%%%%%%%%%%%%%%%%%%%%%%%%%%%%%%%%%%%%%%%%%%%
\zadanie[]
Zdefiniujmy zredukowaną funkcję rozkładu, która mierzy gęstość 
prawdopodobieństwa znalezienia układu w punkcie przestrzeni fazowej 
o danej wartości współrzędnej $x$ (bez względu na wartość $y$):
\begin{align*}
W(x) = \int_0^1 \rho(x, y)dy
\end{align*}

Taki, wielowymiarowy rozkład, odcałkowany po wybranych współrzędnych,
nazywany też jest rozkładem brzegowym, lub zmarginalizowanym. Proszę:
\begin{itemize}
    \item pokazać, że $W_n(x)$, gdzie $n$ oznacza liczbę iteracji przekształcenia 
    piekarskiego spełnia następujące dyskretne równanie Boltzmanna:
\begin{align*}
W_n(x) = \frac{1}{2}\left[W_{n-1}\left(\frac{x}{2}\right) + W_{n-1}\left(\frac{x+1}{2}\right)\right]
\end{align*}
\item znaleźć rozkład równowagowy $W_{\text{eq}}(x)$, czyli taki, który 
      jest niezmienniczy względem przekształcenia
\item pokazać, że entropia Gibbsa: $S_G = -k\int_0^1 W(x)\log W(x)dx$ 
      nie maleje podczas kolejnych iteracji przekształcenia
\end{itemize}

\begin{solution}

Gęstość przestrzeni fazowej w danej iteracji jest dana przez:
\begin{align*}
\rho_n(x, y) = \rho_{n-1}(\hat{B}^{-1}(x, y)) =
 \begin{cases} \rho_{n-1}(x/2, 2y), & y < 1/2 \\ \rho_{n-1}((x+1)/2, 2y - 1), & y \geq 1/2 \end{cases}
\end{align*} 

Całkujemy po $y$ i zamieniamy zmienne $y \to z = 2y$ w pierwszej całce 
oraz $y \to z = 2y - 1$ w drugiej całce:
\begin{align*}
W_n(x) = \int_0^1 \rho_n(x, y) dy = \int_0^{1/2} \rho_{n-1}(x/2, 2y) dy + \int_{1/2}^1 \rho_{n-1}((x+1)/2, 2y - 1) dy = \\
\frac{1}{2}\int_0^1 \rho_{n-1}(x/2, z) dz + \frac{1}{2}\int_0^1 \rho_{n-1}((x+1)/2, z) dz = \\
\frac{1}{2}\left[W_{n-1}\left(\frac{x}{2}\right) + W_{n-1}\left(\frac{x+1}{2}\right)\right]
\end{align*}

Kolejne iteracje przekształcenia mieszają rozkład $W(x)$, dla różnych
wartości $x$ biorąc średnią z wartości $W$ w punktach $x/2$ i $(x+1)/2$. 
Rozkład równowagowy, niezmienniczy względem tego
przekształcenia, to rozkład jednorodny: $W_{\text{eq}}(x) = 1$.

Jaka jest ogólna postać rozwiązania równania Boltzmanna? Możemy rozwinąć $W(x)$ 
w szereg Fouriera:
\begin{align*}
W_{n}(x) = \sum_{k=-\infty}^{\infty} c_k(n) e^{2\pi i k x}
\end{align*}

Wstawiamy to do równania Boltzmanna i korzystamy z ortogonalności funkcji wykładniczych:
\begin{align*}
\sum_{k=-\infty}^{\infty} c_k(n) e^{2\pi i k x} = \frac{1}{2}\left[ \sum_{k=-\infty}^{\infty} c_k(n-1) e^{2\pi i k x/2} + \sum_{k=-\infty}^{\infty} c_k(n-1) e^{2\pi i k (x+1)/2} \right] = \\
\sum_{k=-\infty}^{\infty} c_k(n-1) e^{2\pi i k x/2} \frac{1 + e^{2\pi i k / 2}}{2} = 
\sum_{k=-\infty}^{\infty} c_k(n-1) e^{2\pi i k x/2} \frac{1 + (-1)^k}{2} = \\
\sum_{k=-\infty}^{\infty} c_k(n-1) e^{2\pi i k x/2} \delta_{k 2j} = 
\sum_{j=-\infty}^{\infty} c_{2j}(n-1) e^{2\pi i j x}
\end{align*}

\begin{align*}
\sum_{k=-\infty}^{\infty} c_k(n+1) e^{2\pi i k x} = 
\frac{1}{2}\left[ \sum_{j=-\infty}^{\infty} c_{2j}(n-1) e^{2\pi i j x/2} + 
\sum_{j=-\infty}^{\infty} c_{2j}(n-1) e^{2\pi i j (x+1)/2} \right] = \\
\sum_{j=-\infty}^{\infty} c_{2j}(n-1) e^{2\pi i j x/2} \frac{1 + e^{2\pi i j / 2}}{2} = 
\sum_{j=-\infty}^{\infty} c_{2j}(n-1) e^{2\pi i j x/2} \delta_{j 2l} 
\end{align*}

Wniosek: w każdej iteracji przekształcenia, 
współczynniki $c_k$ dla nieparzystych $k$ (w danej iteracji) stają się zerowe, 
a współczynniki dla parzystych $k$ pozostają bez zmian. 
Ostatecznie, po nieskończenie wielu iteracjach, pozostaje 
tylko współczynnik $c_0$, który odpowiada rozkładowi jednorodnemu.

\vspace{1cm}

Ewolucję entropii Gibbsa znajdujemy korzystając z postaci $W_n(x)$:
\begin{align*}
S_n = -k\int_0^1 W_n(x)\log W_n(x) dx = \\
-k\int_0^1 \frac{1}{2}\left[W_{n-1}\left(\frac{x}{2}\right) + W_{n-1}\left(\frac{x+1}{2}\right)\right] \log\left( \frac{1}{2}\left[W_{n-1}\left(\frac{x}{2}\right) + W_{n-1}\left(\frac{x+1}{2}\right)\right] \right) dx
\end{align*}

Nie obliczamy jawnej postaci $S(n)$, ale chcemy jedynie zbadać 
czy $S(n) \geq S(n-1)$. Funkcja podcałkowa ma postać:
\begin{align*}
\frac{1}{2}\left[W_{n-1}\left(\frac{x}{2}\right) + W_{n-1}\left(\frac{x+1}{2}\right)\right] \log\left( \frac{1}{2}\left[W_{n-1}\left(\frac{x}{2}\right) + W_{n-1}\left(\frac{x+1}{2}\right)\right] \right) = \\
f(a, b) = \frac{a + b}{2} \log\left( \frac{a + b}{2} \right) = f(z)\log(z), \text{ gdzie } z = \frac{a + b}{2}
\end{align*}

Funkcja $f(z) = z\log(z)$ jest wypukła: druga pochodna jest dodatnia dla $z > 0$:
\begin{align*}
f''(z) = \frac{1}{z} > 0
\end{align*}

Wypukłość funkcji oznacza, że prosta łącząca 
dwa punkty na wykresie funkcji leży powyżej wykresu funkcji:
\begin{align*}
  f''(z) \simeq \frac{f(z+h) - 2f(z) + f(z-h)}{(h)^2} > 0 \Rightarrow \\
f(z) < \frac{f(a) + f(b)}{2} \text{ dla } z = \frac{a + b}{2}
\end{align*}

Fakt ten jest znany jako nierówność Jensena. W naszym przypadku:
\begin{align*}
f\left(\frac{a + b}{2}\right) \leq \frac{f(a) + f(b)}{2} \Rightarrow \\
\frac{a + b}{2} \log\left( \frac{a + b}{2} \right) \leq \frac{a\log(a) + b\log(b)}{2}
\end{align*}

Zatem:
\begin{align*}
S_n = -k\int_0^1 f(a, b) dx \geq -k\int_0^1 \frac{a\log(a) + b\log(b)}{2} dx = \\
-\frac{k}{2} \int_0^1 W_{n-1}\left(\frac{x}{2}\right) \log\left(W_{n-1}\left(\frac{x}{2}\right)\right) dx - 
\frac{k}{2} \int_0^1 W_{n-1}\left(\frac{x+1}{2}\right) \log\left(W_{n-1}\left(\frac{x+1}{2}\right)\right) dx = \\
|| z = \frac{x}{2} \text{ w pierwszej całce, } z = \frac{x+1}{2} \text{ w drugiej całce} || = \\
-k\int_0^{1/2} W_{n-1}(z) \log\left(W_{n-1}(z)\right) dz -k \int_{1/2}^1 W_{n-1}(z) \log\left(W_{n-1}(z)\right) dz = \\
-k\int_0^1 W_{n-1}(z) \log\left(W_{n-1}(z)\right) dz =
S_{n-1}
\end{align*}

Zatem $S_n \geq S_{n-1}$, czyli entropia Gibbsa nie maleje podczas kolejnych iteracji 
przekształcenia piekarskiego.


\newpage
\end{solution}
%%%%%%%%%%%%%%%%%%%%%%%%%%%%%%%%%%%%%%%%%%%%%%%%%%%%%%%%%%%%%%%%%%%%%%%
%%%%%%%%%%%%%%%%%%%%%%%%%%%%%%%%%%%%%%%%%%%%%%%%%%%%%%%%%%%%%%%%%%%%%%%%%%%%%%%%
\end{document}